\subsection{EXAMPLES: SIMPLICIAL CONES AND SIMPLICES}

\begin{enumerate}
\def\labelenumi{\alph{enumi})}
\tightlist
\item
  Let \(a\) be a point of \(\mathbf{E}\) , and \((e_1, \ldots, e_d)\) a basis of \(\mathbf{T}\) ; every point of \(\mathbf{E}\) can thus be written uniquely in the form
\end{enumerate}

\[
x = a + \xi_ {1}. e _ {1} + \dots + \xi_ {d}. e _ {d}\tag{7}
\]

with \(\xi_1, \ldots, \xi_d\) being real numbers. Denote by \(e_i'\) the affine function on E which, at any \(x \in \mathbf{E}\) , written in the form (7), takes the value \(\xi_i\) . Denote by \(\mathbf{H}_i\) the hyperplane formed by those \(x\) such that \(e_i'(x) = 0\) , and by \(\mathfrak{H}\) the set of hyperplanes \(\mathbf{H}_1, \ldots, \mathbf{H}_d\) . For every subset J of the set I = \{1, 2, \ldots, d\}, put \(\mathbf{H}_J = \bigcap_{i \in J} \mathbf{H}_i\) ; for every sequence \((\varepsilon_1, \ldots, \varepsilon_d)\) of numbers equal to 0, 1 or -1, denote by \(\mathbf{F}(\varepsilon_1, \ldots, \varepsilon_d)\) the set of those \(x \in \mathbf{E}\) such that \(e_i'(x)\) is of sign \(\varepsilon_i\) for all \(i\) in I (General Topology, Chap. 4, §3, no. 2). It is immediate that the facets defined by \(\mathfrak{H}\) in E are the sets \(\mathbf{F}(\varepsilon_1, \ldots, \varepsilon_d)\) and that these sets are pairwise distinct; the support of \(\mathbf{F}(\varepsilon_1, \ldots, \varepsilon_d)\) is equal to \(\mathbf{H}_J\) if J is the set of \(i \in I\) such that \(\varepsilon_i = 0\) ; in particular, the chambers are the sets of the form \(\mathbf{F}(\varepsilon_1, \ldots, \varepsilon_d)\) where each of the numbers \(\varepsilon_i\) is equal to 1 or -1.

The set \(\mathrm{C} = \mathrm{F}(1, \ldots, 1)\) formed by those \(x\) with \(e_i'(x) > 0\) for all \(i \in \mathrm{I}\) is a chamber, called the open simplicial cone with vertex \(a\) defined by the basis \((e_{1},\ldots,e_{d})\) . Its closure consists of the points x such that \(e_{i}^{\prime}(x)\geqslant0\) for all \(i\in I\) when \(d\geqslant1\) ; otherwise, its closure is empty. For every subset J of I, let \(C_{J}\) be the set of points x of E such that \(e_{i}^{\prime}(x)=0\) for \(i\in J\) and \(e_{i}^{\prime}(x)>0\) for \(i\in I-J\) . Then \(C_{J}\) is a facet with support \(H_{J}\) , and it is an open simplicial cone with vertex a in the affine space \(H_{J}\) ; moreover, \(\overline{C}=\bigcup_{J\subset I}C_{J}\) . In particular, the walls of C are the hyperplanes \(H_{i}\) for \(i\in I\) , and the face of C contained in \(H_{i}\) is equal to \(C_{\{i\}}\) .

None of the sets \(\mathrm{H}_i, \mathrm{H}_J, \mathrm{C}, \mathrm{C}_J\) and \(\mathrm{F}(\varepsilon_1, \ldots, \varepsilon_d)\) change if we replace the basis \((e_1, \ldots, e_d)\) by a basis \((\lambda_1 e_1, \ldots, \lambda_d e_d)\) with \(\lambda_i > 0\) for all \(i\) .

\begin{enumerate}
\def\labelenumi{\alph{enumi})}
\setcounter{enumi}{1}
\tightlist
\item
  Suppose now that we are given an affinely free system of points of E, say \((a_{0}, a_{1}, \ldots, a_{d})\) . We know that every point of E can be written uniquely in the form \(x = \xi_{0}.a_{0} + \cdots + \xi_{d}.a_{d}\) , where \(\xi_{0}, \ldots, \xi_{d}\) are real numbers with \(\xi_{0} + \cdots + \xi_{d} = 1\) (Algebra, Chap. 2, §9, no. 3). Define affine functions \(f_{0}, \ldots, f_{d}\) , the function \(f_{i}\) associating to the point x, written as above, the number \(\xi_{i}\) . Denote by \(H_{i}\) the hyperplane of E defined by \(f_{i}(x) = 0\) and by H the set of hyperplanes \(H_{0}, H_{1}, \ldots, H_{d}\) ; finally, put \(I = \{0, 1, \ldots, d\}\) . The open simplex with vertices \(a_{0}, \ldots, a_{d}\) is the set C of points x of E such that \(f_{i}(x) > 0\) for all \(i \in I\) ; it is one of the chambers defined by H in E. The closure of C is the set \(\overline{C}\) of points x such that \(f_{i}(x) \geqslant 0\) for all \(i \in I\) ; it is the convex envelope of the finite set \(\{a_{0}, \ldots, a_{d}\}\) and it is easy to see that the extreme points of \(\overline{C}\) are \(a_{0}, \ldots, a_{d}\) .
\end{enumerate}

For any subset \(J\) of \(I\) distinct from \(I\) , put \(H_J = \bigcap_{i \in J} H_i\) and let \(C_J\) be the set of points \(x\) of \(E\) such that \(f_i(x) = 0\) for \(i \in J\) and \(f_i(x) > 0\) for \(i \in I - J\) . The set \(C_J\) is an open simplex in the affine space \(H_J\) with vertices the points \(a_i\) for \(i \in I - J\) ; we have \(C_\emptyset = C, \overline{C} = \bigcup_{J \neq I} C_J\) and \(C_J \neq C_{J'}\) if \(J \neq J'\) ; moreover, \(C_J\) is a facet with support \(H_J\) . In particular, the walls of \(C\) are \(H_0, \ldots, H_d\) and \(C_{\{i\}}\) is the face contained in \(H_i\) .

For any non-empty subset \(K\) of \(I\) , let \(B_K\) be the set of points \(x\) of \(E\) such that \(f_i(x) > 0\) for \(i \in K\) and \(f_i(x) < 0\) for \(i \in I - K\) . The sets \(B_K\) are the chambers defined by \(\mathfrak{H}\) in \(E\) and we have \(B_I = C\) . It is easy to see that \(\overline{C}\) is compact; on the other hand, if \(K\) is a subset of I distinct from \(I\) , of cardinal \(p\) , the chamber \(B_K\) contains the sequence of points \(x_n\) defined for \(n \geqslant 2\) by

\[
f _ {i} (x _ {n}) = \left\{ \begin{array}{l l} n & \text { for } i \in K \\ (1 - p n) / (d + 1 - p) & \text { for } i \in I - K \end{array} \right.
\]

showing that \(\mathsf{B}_{\mathsf{K}}\) is not relatively compact.

